\documentclass[11pt]{article}

\usepackage[margin=1in]{geometry}
\usepackage{amsmath,amssymb,amsthm}
\usepackage{booktabs}
\usepackage{array}
\usepackage{longtable}
\usepackage{hyperref}

\newtheorem{theorem}{Theorem}
\newtheorem{lemma}{Lemma}

\newcommand{\YI}{Y_I}
\newcommand{\YS}[1]{Y_{S_{#1}}}
\newcommand{\U}[2]{U_{#1#2}}
\newcommand{\IO}{\operatorname{IO}}

\title{A Detailed Stage 4 Proof for the Okamura 4N Network}
\author{Generated from the RL Stage 4 proof trace}
\date{June 16, 2026}

\begin{document}
\maketitle

\section{Network, sessions, and target statement}

The Okamura 4N instance has node set
\[
V=\{a,b,c,d\},
\]
edge set
\[
E=\{ab,bc,cd,ad,ac\},
\]
and three unicast sessions
\[
I=\{(a,c),(b,d),(a,b)\}.
\]
All edge capacities are normalized to one use of the edge alphabet per network
use.  If each session has rate \(r\), then the joint session entropy satisfies
\[
h(\YI)=3r\log_2 b,
\]
where \(b\) is the alphabet size.  The Stage 4 run used the Phase 1 partition
\[
P_1=\{a\},\qquad P_2=\{c\},\qquad P_3=\{b,d\},
\]
whose partition bound is
\[
r\leq 1.25.
\]
The Stage 4 terminal certificate recorded in
\texttt{config\_files/stage4\_proof\_log.json} and
\texttt{text\_files/proof\_okamura\_4N.txt} is
\begin{equation}
\label{eq:terminal}
8.200\,h(\YI)
\leq
2.240\,h(\U{a}{b})
+6.270\,h(\U{b}{c})
+7.760\,h(\U{c}{d})
+3.730\,h(\U{a}{d})
+4.830\,h(\U{a}{c}).
\end{equation}
This proof explains the node I/O inequalities used to build the certificate,
the fractional Stage 4 trace that produced the recorded edge vector, and the
rate extraction from \eqref{eq:terminal}.

\begin{theorem}
For the Okamura 4N network above, the Stage 4 terminal inequality implies
\[
r\leq 1.009350.
\]
This improves the partition bound \(1.25\) by \(19.252\%\).
\end{theorem}

\section{Node I/O inequalities}

The Phase 3 primitive is the per-node I/O inequality.  For a node \(v\), the
code computes
\[
\frac{\#\{\hbox{sessions touching }v\}}{|I|}h(\YI)
\leq
h(\YS{v})+\sum_{e\ni v}h(U_e).
\]
Here ``touching'' means that \(v\) is either the source or the sink of the
session.  In this graph,
\[
\#\operatorname{touch}(a)=2,\quad
\#\operatorname{touch}(b)=2,\quad
\#\operatorname{touch}(c)=1,\quad
\#\operatorname{touch}(d)=1.
\]
Therefore the four raw node I/O constraints are
\begin{align}
\IO(a):\quad
\frac23 h(\YI)
&\leq h(\YS{a})+h(\U{a}{b})+h(\U{a}{d})+h(\U{a}{c}), \label{eq:ioa}\\
\IO(b):\quad
\frac23 h(\YI)
&\leq h(\YS{b})+h(\U{a}{b})+h(\U{b}{c}), \label{eq:iob}\\
\IO(c):\quad
\frac13 h(\YI)
&\leq h(\YS{c})+h(\U{b}{c})+h(\U{c}{d})+h(\U{a}{c}), \label{eq:ioc}\\
\IO(d):\quad
\frac13 h(\YI)
&\leq h(\YS{d})+h(\U{c}{d})+h(\U{a}{d}). \label{eq:iod}
\end{align}
Each inequality is a standard encoding inequality: the outgoing variables at
node \(v\) are deterministic functions of the incoming variables and the source
symbols available at \(v\), so the entropy carried by the node output cannot
exceed the entropy of the node input.

\section{Fractional I/O rule}

Stage 4 extends the integer partition proof family by forming convex
combinations of node I/O inequalities:
\[
F(u,v;\lambda)=\lambda \IO(u)+(1-\lambda)\IO(v),
\qquad 0<\lambda<1.
\]
For example, using \eqref{eq:ioc} and \eqref{eq:iod},
\[
F(c,d;0.60)
=0.60\IO(c)+0.40\IO(d),
\]
which expands to
\[
\frac13h(\YI)
\leq
0.60h(\YS{c})+0.40h(\YS{d})
+0.60h(\U{b}{c})
+h(\U{c}{d})
+0.40h(\U{a}{d})
+0.60h(\U{a}{c}).
\]
The logged proof repeatedly constructs such fractional inequalities and places
selected ones into the accumulator.  Addition of valid inequalities remains
valid, so any nonnegative sum of accumulated inequalities is valid.

\section{The logged Stage 4 trace}

The best Okamura 4N episode used for the proof document is episode 2657.  The
logger records \(41\) actions:
\[
26\ \hbox{FRACTIONAL\_IO},\qquad
13\ \hbox{ADD\_TO\_ACCUMULATOR},\qquad
2\ \hbox{STORE\_AND\_RESET}.
\]
No crypto, decode, cross-submodularity, or plain submodularity action is
recorded for this episode.  The following table lists the fractional node I/O
actions.  If the action is \(u,v,\lambda\), it means
\(\lambda\IO(u)+(1-\lambda)\IO(v)\).

\begin{longtable}{cclc}
\toprule
Step & Action & Nodes & \(\lambda\)\\
\midrule
\endhead
1  & FRACTIONAL\_IO & \(c,d\) & 0.60\\
2  & FRACTIONAL\_IO & \(b,c\) & 0.50\\
4  & FRACTIONAL\_IO & \(a,b\) & 0.40\\
5  & FRACTIONAL\_IO & \(a,d\) & 0.75\\
7  & FRACTIONAL\_IO & \(c,d\) & 0.33\\
9  & FRACTIONAL\_IO & \(c,d\) & 0.75\\
10 & FRACTIONAL\_IO & \(c,d\) & 0.67\\
11 & FRACTIONAL\_IO & \(b,c\) & 0.75\\
12 & FRACTIONAL\_IO & \(c,d\) & 0.25\\
13 & FRACTIONAL\_IO & \(b,c\) & 0.33\\
14 & FRACTIONAL\_IO & \(c,d\) & 0.60\\
15 & FRACTIONAL\_IO & \(b,c\) & 0.33\\
17 & FRACTIONAL\_IO & \(c,d\) & 0.67\\
19 & FRACTIONAL\_IO & \(c,d\) & 0.75\\
21 & FRACTIONAL\_IO & \(a,d\) & 0.50\\
22 & FRACTIONAL\_IO & \(c,d\) & 0.50\\
23 & FRACTIONAL\_IO & \(b,c\) & 0.50\\
24 & FRACTIONAL\_IO & \(c,d\) & 0.67\\
25 & FRACTIONAL\_IO & \(c,d\) & 0.40\\
26 & FRACTIONAL\_IO & \(b,c\) & 0.25\\
29 & FRACTIONAL\_IO & \(a,c\) & 0.25\\
30 & FRACTIONAL\_IO & \(b,c\) & 0.25\\
34 & FRACTIONAL\_IO & \(a,c\) & 0.33\\
36 & FRACTIONAL\_IO & \(b,c\) & 0.40\\
37 & FRACTIONAL\_IO & \(c,d\) & 0.25\\
41 & FRACTIONAL\_IO & \(c,d\) & 0.60\\
\bottomrule
\end{longtable}

\section{Accumulator arithmetic and source cancellation}

The decisive logged storage event is step 33.  Immediately before this step,
the accumulator contains ten fractional inequalities.  When summed and passed
through the code's \texttt{cancel\_source\_terms} rule, the trace records
\begin{equation}
\label{eq:stored}
3.200\,h(\YI)
\leq
2.240\,h(\U{a}{b})
+6.270\,h(\U{b}{c})
+7.760\,h(\U{c}{d})
+3.730\,h(\U{a}{d})
+4.830\,h(\U{a}{c}).
\end{equation}
The source-cancellation rule is the following.  Let \(m_v\) be the number of
sessions whose source is \(v\).  Here
\[
m_a=2,\qquad m_b=1,\qquad m_c=m_d=0.
\]
If a summed inequality contains RHS source terms
\(\sum_v \alpha_v h(\YS{v})\), then source independence gives
\[
h(\YS{v})=m_v r\log_2 b
=\frac{m_v}{3}h(\YI).
\]
Thus the source contribution can be moved into the left side in
\(h(\YI)\)-units, and the source terms are removed from the terminal form.
That is exactly the rule implemented by \texttt{cancel\_source\_terms}.

The edge vector in \eqref{eq:stored} is also easy to audit from the node I/O
basis.  It is the incidence vector generated by the nonnegative node weights
\[
x_a=0.400,\qquad x_b=1.840,\qquad x_c=4.430,\qquad x_d=3.330,
\]
since
\begin{align*}
x_a+x_b&=2.240,&
x_b+x_c&=6.270,&
x_c+x_d&=7.760,\\
x_a+x_d&=3.730,&
x_a+x_c&=4.830.
\end{align*}
This check confirms that the RHS edge vector is a nonnegative combination of
the valid node I/O RHS edge incidences.

\section{Terminal certificate}

The terminal inequality selected by the Stage 4 episode summary is
\eqref{eq:terminal}.  It has the same RHS edge vector as the step-33 stored
inequality, but the terminal coefficient of \(h(\YI)\) is \(8.200\).  The
terminal detail recorded by the logger is:
\[
\begin{array}{c|ccccc|c}
\hbox{coefficient of }h(\YI)
& \U{a}{b} & \U{b}{c} & \U{c}{d} & \U{a}{d} & \U{a}{c}
& \hbox{edge total}\\
\hline
8.200 & 2.240 & 6.270 & 7.760 & 3.730 & 4.830 & 24.830.
\end{array}
\]
The proof document treats this terminal line as the final certificate emitted
by the Stage 4 proof logger.  All final rate arithmetic below follows only
from this terminal certificate, capacity normalization, and the session entropy
identity \(h(\YI)=3r\log_2 b\).

\section{Rate extraction}

From \eqref{eq:terminal}, using \(h(U_e)\leq \log_2 b\) for each unit-capacity
edge variable, we get
\[
8.200\,h(\YI)
\leq
(2.240+6.270+7.760+3.730+4.830)\log_2 b.
\]
The RHS coefficient sum is
\[
2.240+6.270+7.760+3.730+4.830=24.830.
\]
Since \(h(\YI)=3r\log_2 b\), cancellation of \(\log_2 b\) gives
\[
8.200\cdot 3r \leq 24.830.
\]
Therefore
\[
r\leq \frac{24.830}{24.600}=1.009349593\ldots,
\]
so the logged Stage 4 bound is
\[
\boxed{r\leq 1.009350}.
\]
Compared with the partition bound \(1.25\), the relative improvement is
\[
\frac{1.25-1.009350}{1.25}\times 100\%
=19.252\%.
\]

\section{Trace audit summary}

\begin{itemize}
\item Training completed on 2026-06-16 at 02:54:16 after 45167.3 seconds.
\item The Okamura 4N entry in \texttt{text\_files/training\_summary.txt}
reports \(PB=1.2500\), \(LP\ LB=1.0000\), \(RL\ Bound=1.009350\), and
\(19.25\%\) improvement.
\item The proof document \texttt{text\_files/proof\_okamura\_4N.txt} records
episode 2657 and the terminal inequality \eqref{eq:terminal}.
\item The Stage 4 JSON logger records the terminal form as valid, with
\(Y_I\)-coefficient \(8.2\), edge total \(24.83\), and no residual source,
\(Y_{ST}\), or \(Y_{I,P_i}\) terms.
\end{itemize}

\end{document}
